\documentclass[twocolumn]{article}
\usepackage[utf8]{inputenc}
\usepackage{amsmath, amssymb, amsthm}
\usepackage{booktabs}
\usepackage{graphicx}
\usepackage{hyperref}
\usepackage{geometry}
\geometry{letterpaper, margin=0.75in}

\title{\textbf{KudaHitamQuant: Ultra-High Fidelity 1-bit KV Cache Compression via Structured Spectral Projections}}
\author{Oki Abrian \\ \small{Independent Researcher}}
\date{March 2026}

\begin{document}

\maketitle

\begin{abstract}
Large Language Models (LLMs) with long context windows are critically bottlenecked by the Key-Value (KV) cache. While Multi-Head Latent Attention (MLA) reduces the latent dimension, the linear growth of the cache prohibits extreme sequence lengths. We introduce \textbf{KudaHitamQuant}, a high-fidelity quantization framework optimized for Qwen-3.5 MLA. By leveraging \textbf{Fast Walsh-Hadamard Transform (FWHT)} and \textbf{Optimal Lloyd-Max Centroids}, we achieve superior reconstruction fidelity through perfectly orthogonal structured transformations. We implement \textbf{Gila Mode}, a raw CUDA engine utilizing \textbf{Warp-Shuffle} primitives for $O(D \log D)$ projection, and an **Asymmetric Attention Score Recovery** mechanism that keeps queries in high precision. Benchmarks on \textbf{Qwen-3.5 2B} demonstrate that KudaHitamQuant achieves **near-lossless fidelity (0.9993)** at **2-bits**, significantly outperforming Gaussian QJL baselines.
\end{abstract}

\section{Introduction}
The quest for long context length has shifted the primary optimization target from compute to KV cache bandwidth. Recent works like Google's TurboQuant explored Quantized Johnson-Lindenstrauss (QJL) but rely on dense Gaussian projections ($O(D^2)$). We propose a structured approach that leverages the spectral properties of the latent space in Qwen-3.5 ($D=256$).

\section{Methodology}
\subsection{Structured Spectral Projection}
KudaHitamQuant projects the KV state $x$ into a spectral domain via:
\begin{equation}
y = H(x \odot d) \cdot \frac{1}{\sqrt{D}}
\end{equation}
where $H$ is the Hadamard matrix and $d$ is a Rademacher random diagonal. Quantization follows via an optimal Lloyd-Max codebook $\mathcal{C} = \{c_1, \dots, c_{2^B}\}$.

\subsection{Asymmetric Score Recovery}
To maintain accuracy, we keep Queries $q$ in high precision while compressing Keys $k$. The attention score is recovered:
\begin{equation}
\text{Score} = (q \odot d)^T \cdot H^T \cdot \text{lookup}(\mathcal{C}_k)
\end{equation}
This eliminates the "noisy" sign-only projection of standard 1-bit QJL.

\subsection{Gila Mode: CUDA Warp-Shuffles}
We implement the $O(D \log D)$ FWHT using CUDA's \texttt{\_\_shfl\_xor\_sync} to eliminate shared memory bank conflicts and global memory roundtrips. Together with a **Global Centroid Cache** and **Lazy JIT Loading**, the engine achieves sub-millisecond latency for $D=256$.

\section{Experimental Analysis}
\begin{table*}[t]
\centering
\caption{Benchmark Results on NVIDIA T4 (Qwen-3.5 2B MLA, $D=256$)}
\begin{tabular}{@{}lcccccc@{}}
\toprule
Context & Task & Mode & Acc (V2/F: KudaHitam) & Acc (G: Gaussian) & Comp(V2/F) & Comp(G) \\ \midrule
10k & Reasoning & 1-bit & \textbf{0.9977} (+0.46\%) & 0.9931 & 0.86ms & 0.49ms \\
10k & Reasoning & 2-bit & \textbf{0.9993} (+0.13\%) & 0.9980 & 0.86ms & 0.44ms \\
10k & Math & 1-bit & \textbf{0.9977} (+0.46\%) & 0.9931 & 0.87ms & 0.47ms \\ \cmidrule{1-7}
40k & Reasoning & 1-bit & \textbf{0.9973} (+0.56\%) & 0.9917 & 0.89ms | 0.49ms \\
40k & Reasoning & 2-bit & \textbf{0.9992} (+0.15\%) & 0.9977 & 0.84ms | 0.47ms \\
40k & Coding & 2-bit & \textbf{0.9991} (+0.16\%) & 0.9975 & 0.97ms | 0.45ms \\ \bottomrule
\end{tabular}
\end{table*}

\section{Conclusion}
\textbf{KudaHitamQuant} sets a new state-of-the-art for KV compression by combining structured orthogonal projections with asymmetric attention recovery, delivering near-lossless 2-bit performance.

\end{document}
